\section{Analysis of Model-\texorpdfstring{$B$}{B} (\texorpdfstring{$\gamma_1,\gamma_2>0$}{gamma1,gamma2>0}, \texorpdfstring{$\theta_1=\theta_2=0$}{theta1=theta2=0})}\label{sec:model_b}

Model-$B$ introduces bidirectional jockeying ($\gamma_1,\gamma_2>0$) while retaining no abandonment ($\theta_1=\theta_2=0$). Jockeying conserves the total number of customers in the system, so the stationary distribution still exists under the same condition as Model-$A$. However, jockeying breaks the per-class Poisson/renewal structure: a class-2 customer can no longer identify its effective service time with a class-1 busy period, so the Pollaczek--Khinchine probabilistic route is unavailable. Applying the method of characteristics to the fundamental equation yields a general solution, but closing it~---~i.e., determining the unknown boundary function $P_y(y)$~---~requires solving a Volterra-type integral equation that does not admit a closed form. Model-$B$ is therefore intentionally left open, in structural parallel with Model-$X$; its role is to motivate the tractable specialisation Model-$B_2$.

\textbf{Stability.} $\rho = (\lambda_1+\lambda_2)/\mu < 1$. Jockeying conserves the total customer count and does not affect this condition. Figure~\ref{fig:model-b-queue} shows the queue layout.

\begin{figure}[H]
    \centering
    \resizebox{0.62\textwidth}{!}{%
        % Model B: θ₁=θ₂=0, γ₁,γ₂>0
% Bidirectional jockeying between queues (dashed). No abandonment.
\begin{tikzpicture}[>=Latex, thick]

    % ── Queue 1 (top, Class-1, priority) ──────────────────────────────────────
    \draw (0, 1.9) -- (2, 1.9) -- (2, 1.1) -- (0, 1.1);
    \foreach \x in {0.2, 0.4, ..., 1.8} { \draw[thin] (\x, 1.8) -- (\x, 1.2); }
    \draw[->] (-1.5, 1.5) -- (-0.1, 1.5) node[midway, above] {$\lambda_1$};

    % ── Queue 2 (bottom, Class-2) ──────────────────────────────────────────────
    \draw (0, -1.1) -- (2, -1.1) -- (2, -1.9) -- (0, -1.9);
    \foreach \x in {0.2, 0.4, ..., 1.8} { \draw[thin] (\x, -1.2) -- (\x, -1.8); }
    \draw[->] (-1.5, -1.5) -- (-0.1, -1.5) node[midway, above] {$\lambda_2$};

    % ── Jockeying between queues ──────────────────────────────────────────────
    % γ₁: Class-1 jockeys down to Class-2 queue
    \draw[->, dashed] (0.6, 1.0) -- (0.6, -1.0) node[midway, left] {$\gamma_1$};
    % γ₂: Class-2 jockeys up to Class-1 queue
    \draw[->, dashed] (1.4, -1.0) -- (1.4, 1.0) node[midway, right] {$\gamma_2$};

    % ── Server ────────────────────────────────────────────────────────────────
    \node[draw, circle, minimum size=1.2cm] (server) at (5, 0) {$\mu$};
    \draw[->] (2.2,  1.5) -- (server);
    \draw[->] (2.2, -1.5) -- (server);
    \draw[->] (server) -- (7, 0);

\end{tikzpicture}
%
    }
    \caption{Queue layout for Model-$B$. Dashed arrows indicate bidirectional
    jockeying: $\gamma_1$ moves a class-1 customer down to Queue~2 and $\gamma_2$
    moves a class-2 customer up to Queue~1. Both transfers conserve the total~$n$.}
    \label{fig:model-b-queue}
\end{figure}

\begin{corollary}\label{cor:B:pi00}
    The empty-system probability $\pi_0$ (server idle, no one waiting) and the
    both-queues-empty probability $\pi(0,0)$ (server busy, $N_1=N_2=0$) of Model-$B$
    coincide with the Model-$A$ baseline,
    \begin{equation}
        \pi_0 = 1-\rho, \qquad \pi(0,0) = \rho(1-\rho),
        \label{eq:B:pi00}
    \end{equation}
    and the diagonal of the PGF is $P(z,z)=\pi(0,0)/(1-\rho z)$.
\end{corollary}

\begin{proof}
    Setting $x=y$ in the fundamental equation~\eqref{eq:gen:fundamental_equation} with
    $\theta_1=\theta_2=0$, both jockeying terms carry the factor $(x-y)$ and vanish, as
    does the boundary term $\mu(x-y)P_y(y)$; dividing the surviving terms by the common
    factor $x$ leaves
    \begin{equation*}
        \bigl[(\lambda_1+\lambda_2+\mu)x - \mu - (\lambda_1+\lambda_2)x^2\bigr]P(x,x)
        = -\mu(1-x)\,\pi(0,0).
    \end{equation*}
    Writing $\lambda=\lambda_1+\lambda_2$, the bracket factors,
    \begin{equation}
        (\lambda+\mu)x-\mu-\lambda x^{2}
        = -\bigl(\lambda x^{2}-(\lambda+\mu)x+\mu\bigr)
        = -(\lambda x-\mu)(x-1)
        = (\mu-\lambda x)(x-1);
        \nonumber
    \end{equation}
    cancelling $(x-1)$ against $-(1-x)=(x-1)$ on the right gives
    \begin{equation}
        P(x,x) = \frac{\mu\,\pi(0,0)}{\mu-\lambda x} = \frac{\pi(0,0)}{1-\rho x}.
        \label{eq:B:Pxx}
    \end{equation}
    The idle state $(0)$ exchanges probability flux only with the both-queues-empty busy
    state, so the global balance across the idle--busy cut reads
    $\lambda\pi_0=\mu\pi(0,0)$ (an arrival at total rate $\lambda$ leaves $(0)$; a service
    completion at rate $\mu$ returns to it), whence $\pi(0,0)=\rho\pi_0$. Evaluating
    \eqref{eq:B:Pxx} at $x=1$ with $P(1,1)=1-\pi_0$ gives
    $1-\pi_0=\rho\pi_0/(1-\rho)$, so $(1-\pi_0)(1-\rho)=\rho\pi_0$, i.e.\ $\pi_0=1-\rho$
    and $\pi(0,0)=\rho(1-\rho)$. This matches Model-$A$~\eqref{eq:A:pi0_pi00}: jockeying
    conserves the total-customer process and cannot change the marginal idle probability.
    The same identities hold in Models~$B_2$ and~$C_2$ for the same conservation
    reason~\eqref{eq:gen:pi0_noaband}.
\end{proof}

\begin{claim}\label{clm:B:volterra}
    Under stability, $\rho<1$, the bivariate PGF of Model-$B$ satisfies, on each
    characteristic curve $\gamma_2 x + \gamma_1 y = C_1$ of the fundamental equation,
    the integral representation
    \begin{equation}
        P(x,y(x)) \;=\; P_h(x,y(x))\int_{0}^{x}
        \frac{h\bigl(t;\,C_1\bigr)}{P_h\bigl(t,\,y(t)\bigr)}\,dt,
        \label{eq:B:integral_rep}
    \end{equation}
    where $P_h$ is the homogeneous solution~\eqref{eq:B:Ph}, $y(t)=(C_1-\gamma_2 t)/\gamma_1$
    is the characteristic parametrisation, and $h(t;C_1)$ is an explicit forcing
    function depending on the unknown boundary trace $P_y(y)=P(0,y)$.
    Evaluating~\eqref{eq:B:integral_rep} along the diagonal $x=y=z$ and using the
    conservation identity $P(z,z)=\pi(0,0)/(1-\rho z)$ reduces the problem to a
    \emph{Volterra integral equation of the first kind} for $P_y$:
    \begin{equation}
        \mathcal{F}(z) \;=\; \int_0^z \mathcal{K}(z,t)\,P_y\!\bigl(y(t;z)\bigr)\,dt,
        \label{eq:B:volterra}
    \end{equation}
    where $\mathcal{K}$ and $\mathcal{F}$ are fully explicit
    (defined in~\eqref{eq:B:kernel}--\eqref{eq:B:rhs} below). The kernel
    $\mathcal{K}$ does not factorise for general $(\gamma_1,\gamma_2)$,
    so~\eqref{eq:B:volterra} does not reduce to a closed-form solution for $P_y$.
    Model-$B$ is therefore left open.
\end{claim}

\begin{proof}[Proof of Claim~\ref{clm:B:volterra}]
The proof proceeds in five steps: (i)~homogeneous solution via characteristics,
(ii)~rigorous sign analysis of the exponent $E$, (iii)~particular solution via
variation of parameters, (iv)~pinning the free function from the boundary condition,
and (v)~diagonal evaluation yielding the Volterra equation.

% ─────────────────────────────────────────────────────────────────────────────
\medskip\noindent\textbf{Step~1: Homogeneous solution via the method of characteristics.}

The fundamental equation~\eqref{eq:gen:fundamental_equation} with
$\theta_1=\theta_2=0$ is a first-order linear PDE in $P(x,y)$. Setting the
right-hand side to zero and dividing through by $y$ yields the homogeneous PDE
\begin{equation*}
    \gamma_1 x(x-y)\,\frac{\partial P}{\partial x}
    - \gamma_2 x(x-y)\,\frac{\partial P}{\partial y}
    = -\bigl[(\lambda_1+\lambda_2+\mu)x - \mu - \lambda_1 x^2
             - \lambda_2 xy\bigr]P.
\end{equation*}
Following the method of characteristics (\S\ref{sssec:moc}), we associate to
this PDE the characteristic system
\begin{equation*}
    \frac{dx}{\gamma_1 x(x-y)}
    = \frac{dy}{-\gamma_2 x(x-y)}
    = \frac{dP}{-\bigl[(\lambda_1+\lambda_2+\mu)x - \mu
               - \lambda_1 x^2 - \lambda_2 xy\bigr]P}.
\end{equation*}
Taking the ratio of the first two differentials, the $x(x-y)$ factors cancel,
giving $dx/\gamma_1 = dy/(-\gamma_2)$.  Integrating yields the first integral
\begin{equation*}
    C_1 = \gamma_2 x + \gamma_1 y \quad\text{(constant along each characteristic)}.
\end{equation*}
The characteristics are therefore the straight lines of slope $-\gamma_2/\gamma_1$
in the $(x,y)$-plane. Fix one such line and parametrise $y$ by
$y(x) = (C_1-\gamma_2 x)/\gamma_1$.  To extract the equation for $P$ along this
curve we need the coefficient $A(x,y)$ and the denominator $x(x-y)$ after
substitution.  Setting
\begin{equation}
    A(x,y) := (\lambda_1+\lambda_2+\mu)x - \mu - \lambda_1 x^2 - \lambda_2 xy,
    \label{eq:B:Adef}
\end{equation}
and substituting $y = (C_1-\gamma_2 x)/\gamma_1$ so that
$\lambda_2 xy = (\lambda_2 C_1/\gamma_1)x - (\lambda_2\gamma_2/\gamma_1)x^2$,
we obtain, collecting by powers of $x$,
\begin{align*}
    \begin{split}
    A\bigl(x,y(x)\bigr)
    &= (\lambda_1+\lambda_2+\mu)x - \mu - \lambda_1 x^2
       - \frac{\lambda_2 C_1}{\gamma_1}\,x + \frac{\lambda_2\gamma_2}{\gamma_1}\,x^2 \\
    &= -\frac{\gamma_1\lambda_1-\lambda_2\gamma_2}{\gamma_1}\,x^2
       + \frac{\gamma_1(\lambda_1+\lambda_2+\mu)-\lambda_2 C_1}{\gamma_1}\,x
       - \mu.
    \end{split}
\end{align*}
Similarly, $x - y(x) = [(\gamma_1+\gamma_2)x-C_1]/\gamma_1$, so
$\gamma_1 x(x-y(x)) = x[(\gamma_1+\gamma_2)x-C_1]$.  The $dP/P$ equation from
the characteristic system therefore becomes
\begin{align*}
    \frac{dP}{P}
    &= \frac{-A(x,y(x))}{\gamma_1 x(x-y(x))}\,dx
     = \frac{(\gamma_1\lambda_1 - \lambda_2\gamma_2)x^2
             - \bigl[\gamma_1(\lambda_1+\lambda_2+\mu) - \lambda_2 C_1\bigr]x
             + \gamma_1\mu}
            {\gamma_1 x\bigl[(\gamma_1+\gamma_2)x - C_1\bigr]}\,dx.
\end{align*}
To integrate this, we decompose the integrand via partial fractions.  Writing the
integrand as $(1/\gamma_1)\,I(x)\,dx$ with
$I(x) = K + M/x + N(C_1)/[(\gamma_1+\gamma_2)x - C_1]$,
the partial fraction identity requires
\[
    K(\gamma_1+\gamma_2)x^2
    + \bigl[M(\gamma_1+\gamma_2)+N-KC_1\bigr]x
    - MC_1
    = (\gamma_1\lambda_1-\lambda_2\gamma_2)x^2
      - \bigl[\gamma_1(\lambda_1+\lambda_2+\mu)-\lambda_2 C_1\bigr]x
      + \gamma_1\mu.
\]
Matching the $x^2$, $x^0$, and $x^1$ coefficients in turn:
\begin{equation}
    K = \frac{\gamma_1\lambda_1 - \lambda_2\gamma_2}{\gamma_1+\gamma_2},
    \qquad
    M = -\frac{\gamma_1\mu}{C_1},
    \qquad
    N(C_1) = \gamma_1\!\left[
        \frac{\lambda_1+\lambda_2}{\gamma_1+\gamma_2}C_1
        - (\lambda_1+\lambda_2+\mu)
        + \frac{\mu(\gamma_1+\gamma_2)}{C_1}
    \right].
    \label{eq:B:KMN}
\end{equation}
For $N$, the $x^1$ match gives
$N = -\gamma_1(\lambda_1+\lambda_2+\mu)+\lambda_2 C_1+KC_1-M(\gamma_1+\gamma_2)$;
substituting $K$ and $M$, writing $\lambda=\lambda_1+\lambda_2$,
$\Gamma=\gamma_1+\gamma_2$, and collecting the two $C_1$-proportional terms over the
common denominator $\Gamma$,
\begin{align}
    \begin{split}
        N &= -\gamma_1(\lambda+\mu)
        + C_1\,\frac{\lambda_2\Gamma + \gamma_1\lambda_1-\lambda_2\gamma_2}{\Gamma}
        + \frac{\gamma_1\mu\Gamma}{C_1} \\
        &= -\gamma_1(\lambda+\mu)
        + \frac{\gamma_1\lambda}{\Gamma}\,C_1
        + \frac{\gamma_1\mu\Gamma}{C_1},
    \end{split}
    \nonumber
\end{align}
using $\lambda_2\Gamma+\gamma_1\lambda_1-\lambda_2\gamma_2
=\lambda_2\gamma_1+\gamma_1\lambda_1=\gamma_1\lambda$ in the middle term; factoring
out $\gamma_1$ yields~\eqref{eq:B:KMN}.
Integrating $dP/P = (1/\gamma_1)\,I(x)\,dx$ term by term gives
\begin{equation}
    \ln P = \frac{K}{\gamma_1}x - \frac{\mu}{C_1}\ln x
    + \frac{N(C_1)}{\gamma_1(\gamma_1+\gamma_2)}\ln\bigl[(\gamma_1+\gamma_2)x - C_1\bigr]
    + \ln f(C_1),
    \label{eq:B:lnP}
\end{equation}
where $f(C_1)$ is an arbitrary function of the characteristic constant.
Now observe that
\begin{equation*}
    (\gamma_1+\gamma_2)x - C_1
    = (\gamma_1+\gamma_2)x - (\gamma_2 x + \gamma_1 y)
    = \gamma_1(x-y),
\end{equation*}
so $\ln[(\gamma_1+\gamma_2)x-C_1] = \ln\gamma_1 + \ln(x-y)$. Exponentiating
\eqref{eq:B:lnP} term by term,
\begin{equation}
    P \;=\; f(C_1)\,
    \exp\!\Bigl(\tfrac{K}{\gamma_1}x\Bigr)\cdot
    x^{-\mu/C_1}\cdot
    \gamma_1^{\,N(C_1)/(\gamma_1\Gamma)}\,
    (x-y)^{\,N(C_1)/(\gamma_1\Gamma)};
    \nonumber
\end{equation}
the factor $\gamma_1^{\,N(C_1)/(\gamma_1\Gamma)}$ depends on $(x,y)$ only through
$C_1$ and is absorbed, together with $f$, into a redefined arbitrary function
$F(C_1)$. Reading off the exponent of $(x-y)$:
\begin{equation}
    E(x,y) := \frac{N(C_1)}{\gamma_1(\gamma_1+\gamma_2)}\bigg|_{C_1=\gamma_2 x+\gamma_1 y}
    = \frac{\lambda_1+\lambda_2}{(\gamma_1+\gamma_2)^2}(\gamma_2 x+\gamma_1 y)
    - \frac{\lambda_1+\lambda_2+\mu}{\gamma_1+\gamma_2}
    + \frac{\mu}{\gamma_2 x+\gamma_1 y}.
    \label{eq:B:E}
\end{equation}
Exponentiating~\eqref{eq:B:lnP} yields the general homogeneous solution
\begin{equation}
    P_h(x,y) = F(\gamma_2 x+\gamma_1 y)
    \cdot \exp\!\left(\frac{\gamma_1\lambda_1-\lambda_2\gamma_2}{\gamma_1(\gamma_1+\gamma_2)}x\right)
    \cdot x^{-\mu/(\gamma_2 x+\gamma_1 y)}
    \cdot (x-y)^{E(x,y)},
    \label{eq:B:Ph}
\end{equation}
where $F$ is an arbitrary differentiable function to be determined.

% ─────────────────────────────────────────────────────────────────────────────
\medskip\noindent\textbf{Step~2: The exponent $E(x,y)$ is strictly positive on $(0,1)^2$.}

Recall $\lambda = \lambda_1+\lambda_2$ and $\Gamma = \gamma_1+\gamma_2$. Along any
fixed characteristic, $C_1 = \gamma_2 x + \gamma_1 y$ is constant, so the
exponent~\eqref{eq:B:E} depends only on $u := C_1\in(0,\Gamma)$:
\begin{equation*}
    E(u) = \frac{\lambda}{\Gamma^2}\,u - \frac{\lambda+\mu}{\Gamma} + \frac{\mu}{u}.
\end{equation*}
(For $(x,y)\in(0,1)^2$ we have $u = \gamma_2 x+\gamma_1 y \in (0,\Gamma)$; $E(u)$ is
continuous there, with $E(u)\to+\infty$ as $u\to0^+$ and $E(u)\to0^+$ as
$u\to\Gamma^-$, so it is finite at every interior point although not uniformly bounded.)
Multiplying through by
$u\Gamma^2>0$ gives the equivalent quadratic
\begin{equation*}
    q(u) := \lambda u^2 - (\lambda+\mu)\Gamma u + \mu\Gamma^2.
\end{equation*}
The sign of $E$ on $(0,\Gamma)$ equals the sign of $q(u)$ there.  We
factor $q$ by finding its roots.  The discriminant is
$\Delta = (\lambda+\mu)^2\Gamma^2 - 4\lambda\mu\Gamma^2 = \Gamma^2(\mu-\lambda)^2$,
so (using $\rho<1 \Rightarrow \mu>\lambda \Rightarrow \mu-\lambda>0$)
\begin{equation*}
    u_{1,2} = \frac{(\lambda+\mu)\Gamma \pm \Gamma(\mu-\lambda)}{2\lambda}
    \quad\Longrightarrow\quad
    u_1 = \frac{2\lambda\Gamma}{2\lambda} = \Gamma,
    \quad
    u_2 = \frac{2\mu\Gamma}{2\lambda} = \frac{\Gamma}{\rho} > \Gamma.
\end{equation*}
Hence
\begin{equation*}
    q(u) = \lambda(u - \Gamma)\!\left(u - \tfrac{\Gamma}{\rho}\right).
\end{equation*}
For every $u\in(0,\Gamma)$ both factors are strictly negative
($(u-\Gamma)<0$ and $(u-\Gamma/\rho)<0$ since $\Gamma/\rho>\Gamma$),
so $q(u)=\lambda\cdot(\text{neg})\cdot(\text{neg})>0$, and therefore
\begin{equation*}
    E(x,y) > 0 \quad \text{for all } (x,y)\in(0,1)^2.
\end{equation*}
Note that $E\to0$ as $u\to\Gamma^-$ (i.e.\ as $(x,y)\to(1,1)$), while $E$ is
bounded and bounded away from zero in any compact subset of $(0,1)^2$.

\smallskip\noindent\textit{Sign convention and diagonal limit.}
Because $E>0$, the factor $(x-y)^{E}$ in $P_h$ carries the complex phase
$e^{i\pi E}$ whenever $x<y$. This phase is absorbed into the arbitrary function
$F(C_1)$, and the ratio $P_h(x,y(x))/P_h(t,y(t))$ that appears in the integral
representation is always real and positive: both numerator and denominator lie on
the same characteristic (same $C_1$, same constant value of $E$), so the phases
cancel exactly.  Crucially, since $E>0$ the factor $(y(t;z)-t)^E\to0$ as
$t\to z^-$, meaning $P_h$ formally vanishes on the diagonal $x=y$.  The diagonal
equation that appears in Step~5 below is therefore understood as the
$\lim_{x\to z^-}$ of~\eqref{eq:B:integral_rep_body}; the limit-by-limit analysis
carried out there, using the exact linear relation
$y(t;z)-t=(\Gamma/\gamma_1)(z-t)$, confirms that the $0\cdot\infty$ limit is
finite and equals $\pi(0,0)/(1-\rho z)$.

% ─────────────────────────────────────────────────────────────────────────────
\medskip\noindent\textbf{Step~3: Particular solution via variation of parameters.}

We now work with the full (inhomogeneous) fundamental equation restricted to a
fixed characteristic.  Along the curve $y(x)=(C_1-\gamma_2 x)/\gamma_1$, the
total derivative satisfies $\frac{d}{dx}P(x,y(x)) = \partial_x P + \frac{dy}{dx}\partial_y P = \partial_x P - \frac{\gamma_2}{\gamma_1}\partial_y P$,
so the jockeying terms in~\eqref{eq:gen:fundamental_equation} become
\[
    \gamma_1 xy(x-y)\,\partial_x P + \gamma_2 xy(y-x)\,\partial_y P
    = \gamma_1 xy(x-y)\!\left(\partial_x P - \tfrac{\gamma_2}{\gamma_1}\partial_y P\right)
    = \gamma_1 xy(x-y)\,\frac{d}{dx}P(x,y(x)).
\]
The fundamental equation along the characteristic therefore reads, with its
algebraic coefficient factorised as $y\,A(x,y)$ (cf.~\eqref{eq:B:Adef}),
\begin{align}
    \begin{split}
        y\,A(x,y)\,P + \gamma_1 xy(x-y)\,\frac{dP}{dx} \;=\; G(x,y),
    \end{split}
    \nonumber
\end{align}
and dividing through by $\gamma_1 xy(x-y(x))$
reduces the PDE to the first-order linear ODE (for fixed $C_1$)
\begin{equation}
    \frac{dP}{dx} = f(x;C_1)\,P + h(x;C_1),
    \label{eq:B:ode_char}
\end{equation}
where the coefficient $f$ and the forcing term $h$ are
\begin{align}
    f(x;C_1) &= -\frac{A(x,\,y(x))}{\gamma_1 x\bigl(x-y(x)\bigr)},
    \nonumber\\[4pt]
    h(x;C_1) &= \frac{G(x,\,y(x))}{\gamma_1 x\,y(x)\bigl(x-y(x)\bigr)},
    \label{eq:B:hdef}
\end{align}
with $A$ as in~\eqref{eq:B:Adef} and
\begin{equation}
    G(x,y) := \mu(x-y)\,P_y(y) - \mu x(1-y)\,\pi(0,0).
    \label{eq:B:Gdef}
\end{equation}
Since $P_h(x,y(x))$ satisfies the homogeneous equation
$\frac{d}{dx}P_h = f(x;C_1)\,P_h$, it is the fundamental solution of
ODE~\eqref{eq:B:ode_char}.  Applying variation of parameters (\S\ref{sssec:vop}),
the general solution of the inhomogeneous ODE is
\begin{equation}
    P(x,y(x)) = P_h(x,y(x))\!\left[F(C_1) + \int_{x_0}^{x} \frac{h(t;C_1)}{P_h(t,y(t))}\,dt\right].
    \label{eq:B:variation}
\end{equation}

% ─────────────────────────────────────────────────────────────────────────────
\medskip\noindent\textbf{Step~4: Pinning the free function from the boundary condition at $x=0$.}

Along any characteristic with constant $C_1>0$, the endpoint at $x=0$ is
$y(0)=C_1/\gamma_1$, giving
\[
    P\!\left(0,\,\frac{C_1}{\gamma_1}\right) = P_y\!\left(\frac{C_1}{\gamma_1}\right),
\]
which is finite.  However, $P_h$ contains the factor
$x^{-\mu/C_1}\to+\infty$ as $x\to0^+$ (since $\mu/C_1>0$ for $C_1>0$), while
the integral in~\eqref{eq:B:variation} vanishes as $x_0\to0$ (the integration
interval collapses).  Therefore
\[
    \lim_{x\to0^+}P(x,y(x))
    = \lim_{x\to0^+}P_h(x,y(x))\cdot F(C_1)
    = +\infty\cdot F(C_1),
\]
which can only be finite if $F(C_1)=0$.  Setting $x_0=0$ and $F(C_1)=0$
in~\eqref{eq:B:variation} gives the integral representation of
Claim~\ref{clm:B:volterra}:
\begin{equation}
    P(x,y(x)) = P_h(x,y(x))\int_{0}^{x} \frac{h(t;C_1)}{P_h(t,y(t))}\,dt.
    \label{eq:B:integral_rep_body}
\end{equation}

% ─────────────────────────────────────────────────────────────────────────────
\medskip\noindent\textbf{Step~5: Reduction to the Volterra equation.}

By Corollary~\ref{cor:B:pi00} the empty-state probabilities are $\pi_0=1-\rho$ and
$\pi(0,0)=\rho(1-\rho)$, with diagonal value $P(z,z)=\pi(0,0)/(1-\rho z)$
by~\eqref{eq:B:Pxx}; these are the inputs to the diagonal evaluation below.

We simplify the forcing function $h$ on the characteristic.  Substituting
$G(t,y(t))$ from~\eqref{eq:B:Gdef} into~\eqref{eq:B:hdef}:
\begin{align}
    h(t;C_1)
    &= \frac{\mu(t-y(t))\,P_y(y(t)) - \mu t(1-y(t))\,\pi(0,0)}
             {\gamma_1\,t\,y(t)\,(t-y(t))} \nonumber\\
    &= \frac{\mu\,P_y(y(t))}{\gamma_1\,t\,y(t)}
       - \frac{\mu(1-y(t))\,\pi(0,0)}{\gamma_1\,y(t)\,(t-y(t))}.
    \label{eq:B:h_intermediate}
\end{align}
We will evaluate~\eqref{eq:B:integral_rep_body} along the diagonal characteristic,
which passes through $(z,z)$ and has $C_1=(\gamma_1+\gamma_2)z$.  The parametrisation
along that curve is $y(t;z)=((\gamma_1+\gamma_2)z-\gamma_2 t)/\gamma_1$, so
$y(0;z)=(\gamma_1+\gamma_2)z/\gamma_1>z$ and $y(z;z)=z$.  Therefore $y(t;z)>t$
for all $t\in[0,z)$, which means $t-y(t;z)<0$ throughout the integration interval.
Writing $t-y(t;z)=-(y(t;z)-t)$ in the denominator of~\eqref{eq:B:h_intermediate}:
\begin{equation}
    h\bigl(t;\,(\gamma_1+\gamma_2)z\bigr)
    = \frac{\mu}{\gamma_1\,t\,y(t;z)}\,P_y\!\bigl(y(t;z)\bigr)
    + \frac{\mu\,(1-y(t;z))}{\gamma_1\,y(t;z)\,(y(t;z)-t)}\,\pi(0,0).
    \label{eq:B:h_simplified}
\end{equation}

As established in Step~2, $E>0$, so the homogeneous factor $(x-y)^E$ vanishes on the
diagonal and $P_h\bigl(x,y(x;z)\bigr)\to0$ as $x\to z^-$; simultaneously the integrand
in~\eqref{eq:B:integral_rep_body} develops a non-integrable singularity
$\sim(y(t;z)-t)^{-1-E}$ as $t\to z^-$, so the integral diverges. Their product is the
finite $0\cdot\infty$ limit
\begin{equation*}
    \frac{\pi(0,0)}{1-\rho z}
    = \lim_{x\to z^-} P_h\bigl(x,y(x;z)\bigr)
      \int_0^x \frac{h\bigl(t;(\gamma_1+\gamma_2)z\bigr)}{P_h\bigl(t,y(t;z)\bigr)}\,dt,
\end{equation*}
whose value is $P(z,z)=\pi(0,0)/(1-\rho z)$ by Corollary~\ref{cor:B:pi00}.
Substituting~\eqref{eq:B:h_simplified} and separating the term containing the unknown
$P_y$ from the known remainder, this limit splits as
\begin{align}
    \frac{\pi(0,0)}{1-\rho z}
    &= \lim_{x\to z^-}\frac{\mu\,P_h\bigl(x,y(x;z)\bigr)}{\gamma_1}
       \int_0^x \frac{P_y\bigl(y(t;z)\bigr)}{t\cdot y(t;z)\cdot P_h\bigl(t,y(t;z)\bigr)}\,dt
    \nonumber\\
    &\quad
    + \lim_{x\to z^-}\frac{\mu\,P_h\bigl(x,y(x;z)\bigr)\,\pi(0,0)}{\gamma_1}
       \int_0^x \frac{1-y(t;z)}{y(t;z)\cdot(y(t;z)-t)\cdot P_h\bigl(t,y(t;z)\bigr)}\,dt.
    \label{eq:B:split}
\end{align}
We now evaluate the two limits in~\eqref{eq:B:split} one by one; this both
clarifies in what sense the kernel pairing below is to be read and provides an
independent consistency check of the representation. The geometry on the diagonal
characteristic is exactly linear,
\begin{equation}
    y(t;z)-t \;=\; \frac{(\gamma_1+\gamma_2)z-\gamma_2 t}{\gamma_1}\;-\;t
    \;=\; \frac{\Gamma}{\gamma_1}\,(z-t),
    \label{eq:B:exact_linear}
\end{equation}
and along this characteristic ($C_1=\Gamma z$) the exponent of Step~2 takes the
closed form
\begin{equation*}
    E_z \;:=\; E(\Gamma z)
    \;=\; \frac{\lambda z^{2}-(\lambda+\mu)z+\mu}{\Gamma z}
    \;=\; \frac{(\mu-\lambda z)(1-z)}{\Gamma z}\;>\;0,
    \qquad z\in(0,1).
\end{equation*}
Every factor of $P_h$ in~\eqref{eq:B:Ph} other than the exponential, the power of
$x$, and $(x-y)^{E}$ depends on $(x,y)$ only through $C_1$, so along the
characteristic the homogeneous ratio is \emph{exact}:
\begin{equation}
    \frac{P_h\bigl(x,y(x;z)\bigr)}{P_h\bigl(t,y(t;z)\bigr)}
    \;=\; R(x,t)\,\biggl(\frac{z-x}{z-t}\biggr)^{\!E_z},
    \qquad
    R(x,t):=e^{\frac{K}{\gamma_1}(x-t)}\,\Bigl(\frac{t}{x}\Bigr)^{\!\mu/(\Gamma z)},
    \label{eq:B:ratio_exact}
\end{equation}
with $R$ smooth and $R(z,z)=1$. Substituting~\eqref{eq:B:ratio_exact} into the
\emph{second} limit of~\eqref{eq:B:split} and using~\eqref{eq:B:exact_linear}, its
integrand behaves as $w(z)\,(z-t)^{-1-E_z}$ near $t=z$, with
$w(z)=\mu\,\pi(0,0)\,(1-z)/(\Gamma z)$; the truncated integral diverges like
$(z-x)^{-E_z}/E_z$, the prefactor vanishes like $(z-x)^{E_z}$, and the product
converges:
\begin{align*}
    \begin{split}
        \lim_{x\to z^-}\;(z-x)^{E_z}\!\int^{x}\! w(z)\,(z-t)^{-1-E_z}\,dt
        \;&=\; \frac{w(z)}{E_z}
        \;=\; \frac{\mu\,\pi(0,0)\,(1-z)}{\Gamma z}\cdot
              \frac{\Gamma z}{(\mu-\lambda z)(1-z)} \\
        &=\; \frac{\mu\,\pi(0,0)}{\mu-\lambda z}
        \;=\; \frac{\pi(0,0)}{1-\rho z}\,.
    \end{split}
\end{align*}
The second limit alone thus reproduces the full diagonal value --- an independent
confirmation that the representation~\eqref{eq:B:integral_rep_body} is consistent
with Corollary~\ref{cor:B:pi00}. The \emph{first} limit, by contrast, vanishes: its
integrand behaves as $(z-t)^{-E_z}$ near $t=z$, so for $E_z<1$ the integral
converges and is annihilated by the vanishing prefactor $(z-x)^{E_z}$, while for
$E_z\ge1$ the truncated integral grows only like $(z-x)^{1-E_z}$ (logarithmically
at $E_z=1$) and the product is still $O\bigl((z-x)^{\min(E_z,1)}\ln(z-x)\bigr)\to0$.
Hence \emph{at leading order the limit~\eqref{eq:B:split} reduces to the identity
$\pi(0,0)/(1-\rho z)=\pi(0,0)/(1-\rho z)$ and carries no information about $P_y$}:
the dependence on the boundary function enters only at the next order in $(z-x)$.
Each limit is accordingly a \emph{Hadamard finite part}: writing
$P_h^{\,\mathrm{diag}}(z):=\lim_{x\to z^-}P_h\bigl(x,y(x;z)\bigr)$ (which vanishes
pointwise), the kernel pairing below is the extraction of the
truncation-independent, subleading part of the divergent integrals, and neither
factor is meaningful in isolation.
% TODO(citation): add a standard reference for Hadamard finite-part integrals
% (e.g. Hadamard, Lectures on Cauchy's Problem, 1923) -- the bibliography currently
% lacks a suitable entry; comment only, no PDF output.
Writing the kernel and known term \emph{formally} as
\begin{align}
    \mathcal{K}(z,t) &:= \frac{\mu\cdot P_h^{\,\mathrm{diag}}(z)}
                                {\gamma_1\cdot t\cdot y(t;z)\cdot P_h\bigl(t,y(t;z)\bigr)},
    \label{eq:B:kernel}\\[6pt]
    \mathcal{B}(z) &:= \frac{\mu\cdot P_h^{\,\mathrm{diag}}(z)\cdot\pi(0,0)}{\gamma_1}
                        \int_0^z
                        \frac{1-y(t;z)}{y(t;z)\cdot(y(t;z)-t)\cdot P_h\bigl(t,y(t;z)\bigr)}
                        \,dt,
    \nonumber
\end{align}
\emph{to be read through the regularised limit~\eqref{eq:B:split} rather than pointwise},
and setting $\mathcal{F}(z)=\pi(0,0)/(1-\rho z)-\mathcal{B}(z)$,
equation~\eqref{eq:B:split} becomes the Volterra integral equation of the first kind
\begin{equation}
    \mathcal{F}(z) = \int_0^z \mathcal{K}(z,t)\,P_y\!\bigl(y(t;z)\bigr)\,dt,
    \label{eq:B:rhs}
\end{equation}
confirming Claim~\ref{clm:B:volterra}. The finite-part reading is not a decoration.
As the limit-by-limit analysis shows, the displayed $\mathcal{K}$ and $\mathcal{B}$
are formal --- the prefactor $P_h^{\,\mathrm{diag}}(z)$ vanishes pointwise --- and
the equation for $P_y$ lives entirely at subleading order in $(z-x)$. Giving that
extraction a rigorous meaning, whether by matched asymptotics near $t=z$ or by a
change of variable removing the $(z-t)^{-1-E_z}$ singularity before passing to the
limit, is itself the analytical obstruction: it is this step, together with the
fact that the kernel $\mathcal{K}$ depends on $P_h$ through both arguments and does
not simplify for general $(\gamma_1,\gamma_2)$, that places Model-$B$ beyond the
scope of this thesis. Equation~\eqref{eq:B:rhs} is recorded as the precise form of
the open problem, not solved here; resolving it would require Volterra inversion
methods together with the regularisation just described.
\end{proof}

Model-$B$ is therefore left open at this point.  Its role in the hierarchy is
structural: it identifies the precise obstruction~---~the Volterra equation
for $P_y$~---~that the tractable special cases $B_2$ and $C_2$ circumvent, each
by a different mechanism.  Setting $\gamma_2=0$ collapses the characteristic
family from a two-parameter set of oblique lines to a one-parameter set of
vertical lines, degenerating the PDE to an ODE in $x$ alone; the boundary
function $P_y(y)$ is then pinned by an analyticity argument at the interior
singular point $x=y$.  Model-$C_2$ removes jockeying entirely and pins $P_y$ by
a boundedness argument at the boundary of the domain.  Both specialisations are
analysed in the following sections.
